\centerline{\bf \Huge Домашнее задание}
\centerline{\bf \Huge Вписанный угол}

\begin{flushright}
        \scriptsize{Нет ничего постыдного в том, чтобы упасть!\\ Истинный позор — больше не встать.\\ Синтаро Мидорима. Баскетбол Куроку.}
\end{flushright}

\problem{1} 
a) Точки $A,\ B,\ C\ \text{и}\ D$ лежат на окружности с центром О. Известно, что $\angle AOD = \angle BOC.$
Докажите, что $AD = BC$.

\noindentб) Точки $M,\ N,\ P\ \text{и}\ Q$ лежат на окружности с центром О. Известно, что $MN = PQ$.
Докажите, что $\angle MON = \angle POQ.$

\problem{2} Найдите углы четырёхугольника $ABCD$, вершины которого расположены на окружности, если $\angle ABD = 74^{\circ},\ \angle DBC = 38^{\circ},\ \angle BDC = 65^{\circ}$

\problem{3}  Продолжение высоты $CD$, опущенной из вершины $C$ прямого угла прямоугольного
треугольника $ABC$, делит дугу $AB$ описанной окружности на дуги, одна из которых на 40 градусов
больше другой. Найдите острые углы треугольника.

\problem{4}В треугольнике $АВС$ вписанная окружность касается сторон $АВ$, $ВС$ и $АС$ в точках $K$,
$L$ и $M$ соответственно. Найдите углы треугольника $KLM$, если $\angle A = 48^{\circ},\ \angle B = 72^{\circ}.$

\problem{5}\textbf{\textit{Прямая Симсона}}. Основания перпендикуляров, опущенных на стороны треугольника (или их продолжения) из точки, лежащей на описанной около этого треугольника окружности, лежат на одной прямой.

\problem{6} 
а) В равнобедренном треугольнике $ABC$ ($AB = BC$) проведена биссектриса $AE$. Радиус $OE$
описанной окружности треугольника $AEC$ пересекает биссектрису угла $ACB$ в точке $L$. Докажите, что L лежит на средней линии треугольника $AEC$.

\noindentб) В треугольнике $ABC$ угол $B$ равен $60^{\circ}$ Пусть $AA_1$ и $CC_1$ — биссектрисы этого треугольника. Докажите, что точка, симметричная вершине B относительно прямой $A_1C_1$, лежит на стороне $AC$.